\scanpage{352}
\deleted{\[
\unsure{Z^iC^j}=\Lambda^i\Psi\times\Lambda^{j+1}\Phi
\longrightarrow\unsure{C^{i+j}}=\Lambda^{i+j+1}\Phi
\]}

On \unsure{a alors une loi} de $k$-\unsure{Algèbre graduée}
\unsure{alternée sur} $Z^i(C)=\Lambda^i\Psi$, par le
\unsure{produit extérieur}. \unsure{Question : quand et}
\unsure{comment cela se prolonge-t-il en une}
\unsure{structure de} $\unsure{\frac12|k}$-algèbre différentielle?

\diagram{\begin{tikzcd}
Z^{i+1}\otimes C^j\arrow[r,"\text{donné}"] & Z^{i+j+1}\\
C^i\times Z^j\arrow[u]\arrow[r,"\sigma_{ij}"] & C^{i+j}\arrow[u]\\
Z^i\otimes Z^j\arrow[u]\arrow[r,"\text{donné}"] & Z^{i+j}\arrow[u]\arrow[d]\\
0\arrow[u] & 0
\end{tikzcd}}
\margin{$\big]$}

L'indétermination \unsure{pour trouver} $\sigma_{ij}$ est dans
\[
\operatorname{Hom}_k(Z^{i+1}\otimes Z^j,Z^{i+j})\quad\unsure{c.a.}
\]
\[
\underbrace{\Lambda^{i+1}\Psi\otimes\Lambda^j\Psi}_
{\displaystyle\simeq\Lambda^{i+j+1}\Psi\underset{\swarrow}{\oplus}\ \substack{\text{(quelque chose}\\\text{\unsure{d'homotopie à|d'homologie à}}\ \unsure{a|0}}}
\quad\Lambda^{i+j}\Psi
\]

\mspage{353}
